\documentclass[11pt]{article}

% Plantilla común para material de estudio de Física Matemática II.
% Documento A4, una columna, fuente Computer Modern de 11 puntos
% y márgenes de 2.5 cm por lado.

\usepackage[utf8]{inputenc}
\usepackage[T1]{fontenc}
\usepackage[spanish]{babel}
\usepackage{amsmath,amssymb,mathtools,bm,graphicx,xcolor}
\usepackage[a4paper,margin=2.5cm]{geometry}
\usepackage[
    colorlinks=true,
    linkcolor=red,
    citecolor=green,
    urlcolor=red
]{hyperref}

\spanishdecimal{.}

\setlength{\parindent}{0pt}
\setlength{\parskip}{0.45em}
\setlength{\abovedisplayskip}{6pt}
\setlength{\belowdisplayskip}{6pt}
\setlength{\abovedisplayshortskip}{4pt}
\setlength{\belowdisplayshortskip}{4pt}
\allowdisplaybreaks

% Formato de encabezado y problemas
\newcommand{\encabezado}[3]{%
{\Large\bfseries #1\par}
\vspace{2pt}
{\normalsize #2\hfill #3\par}
\vspace{6pt}\hrule\vspace{8pt}}

\newcommand{\problema}[2]{%
\vspace{0.55em}
\noindent\textbf{#1.}\enspace{\small #2}\par
\vspace{0.3em}}

\newcommand{\inciso}[1]{%
\vspace{0.35em}
\noindent\textbf{(#1)}\enspace}

% Comandos auxiliares
\newcommand{\dd}{\,\mathrm{d}}
\newcommand{\ii}{\mathrm{i}}
\newcommand{\e}{\mathrm{e}}
\newcommand{\sen}{\operatorname{sen}}
\newcommand{\grad}{\nabla}
\newcommand{\laplaciano}{\nabla^2}
\newcommand{\R}{\mathbb{R}}

\begin{document}

\encabezado
{Ecuación de onda 2D en un disco}
{Física Matemática II}
{J.R., G.R.}

\problema{1}{Considere la ecuación de onda bidimensional en un círculo de radio $a$, escrita en coordenadas polares $(\rho,\varphi)$, con borde fijo. Se quiere escribir la solución general en términos de soluciones separables y determinar los coeficientes usando las condiciones iniciales
\[
\psi(\rho,\varphi,0)=\psi_0(\rho,\varphi),
\qquad
\frac{\partial \psi}{\partial t}(\rho,\varphi,0)=\eta(\rho,\varphi).
\]}

\textbf{Solución:} La ecuación a resolver es
\begin{equation}
\frac{\partial^2\psi}{\partial t^2}
=v^2\left[\frac{1}{\rho}\frac{\partial}{\partial\rho}\left(\rho\frac{\partial\psi}{\partial\rho}\right)
+\frac{1}{\rho^2}\frac{\partial^2\psi}{\partial\varphi^2}\right],
\qquad 0<\rho<a.
\label{eq:onda}
\end{equation}
La condición de borde es
\begin{equation}
\psi(a,\varphi,t)=0.
\label{eq:borde}
\end{equation}
Buscamos soluciones separables de la forma $\psi(\rho,\varphi,t)=R(\rho)\Phi(\varphi)T(t)$. Sustituyendo esto en \eqref{eq:onda}, se obtiene
\[
R\Phi T''=v^2\left[\frac{1}{\rho}(\rho R')'\Phi T+\frac{1}{\rho^2}R\Phi''T\right].
\]
Dividiendo por $v^2R\Phi T$ y multiplicando por $\rho^2$, queda
\[
\frac{\rho^2}{v^2}\frac{T''}{T}
=\rho\frac{(\rho R')'}{R}+\frac{\Phi''}{\Phi}.
\]
Separando la parte angular, tomamos
\begin{equation}
\Phi''+m^2\Phi=0,
\qquad m=0,\pm1,\pm2,\ldots,
\label{eq:angular}
\end{equation}
por lo que una base angular compleja es $\Phi(\varphi)=\e^{\ii m\varphi}$. Separando la parte temporal con frecuencia $\omega$, se obtiene
\begin{equation}
T''+\omega^2T=0.
\label{eq:temporal}
\end{equation}
Usando \eqref{eq:angular} y \eqref{eq:temporal} en la ecuación separada, la función radial satisface
\begin{equation}
\rho^2R''+\rho R'+\left(\frac{\omega^2}{v^2}\rho^2-m^2\right)R=0.
\label{eq:radial}
\end{equation}
La ecuación \eqref{eq:radial} se identifica con la ecuación de Bessel haciendo el cambio de variable
\[
x=k\rho,\qquad k=\frac{\omega}{v}.
\]
En efecto, al separar variables las derivadas radiales pasan a actuar solamente sobre $R(\rho)$. Si se escribe $R(\rho)=\widetilde R(x)$, entonces la regla de la cadena da, de forma esquemática,
\[
\frac{\dd R}{\dd \rho}=k\frac{\dd \widetilde R}{\dd x},
\qquad
\frac{\dd^2 R}{\dd \rho^2}=k^2\frac{\dd^2 \widetilde R}{\dd x^2}.
\]
Sustituyendo estas relaciones en \eqref{eq:radial} y usando $\rho=x/k$, todos los factores de $k$ se cancelan y queda la forma estándar de Bessel,
\[
x^2\widetilde R''+x\widetilde R'+(x^2-m^2)\widetilde R=0.
\]
Por ello, la solución regular en $\rho=0$ es
\[
R(\rho)=J_{m}(k\rho),
\qquad k=\frac{\omega}{v}.
\]
Sustituyendo esta solución radial en la condición de borde \eqref{eq:borde}, resulta
\[
J_{m}(ka)=0.
\]
Si $\alpha_{mn}$ es el $n$-ésimo cero positivo de $J_m$, entonces
\begin{equation}
k_{mn}=\frac{\alpha_{mn}}{a},
\qquad
\omega_{mn}=vk_{mn}=\frac{\alpha_{mn}v}{a}.
\label{eq:frecuencia}
\end{equation}
Por lo tanto, las soluciones separables complejas son
\begin{equation}
\psi^{\mathrm{sep}}_{mn\pm}(\rho,\varphi,t)
=J_{m}\left(\frac{\alpha_{mn}\rho}{a}\right)\e^{\ii m\varphi}\e^{\pm\ii\omega_{mn}t},
\label{eq:modos}
\end{equation}
con $m=0,\pm1,\pm2,\ldots$ y $n=1,2,3,\ldots$. Para $m<0$ se usa la relación usual entre funciones de Bessel de orden entero negativo y positivo, por lo que los ceros pueden entenderse en términos de $J_{|m|}$.

Con \eqref{eq:modos}, la solución general se escribe como
\begin{equation}
\psi(\rho,\varphi,t)
=\sum_{m=-\infty}^{\infty}\sum_{n=1}^{\infty}
J_{m}\left(\frac{\alpha_{mn}\rho}{a}\right)\e^{\ii m\varphi}
\left[A_{mn}\e^{\ii\omega_{mn}t}+B_{mn}\e^{-\ii\omega_{mn}t}\right].
\label{eq:general}
\end{equation}

Sustituyendo $t=0$ en \eqref{eq:general} y usando $\psi(\rho,\varphi,0)=\psi_0(\rho,\varphi)$, queda
\begin{equation}
\psi_0(\rho,\varphi)
=\sum_{m,n}J_{m}\left(\frac{\alpha_{mn}\rho}{a}\right)(A_{mn}+B_{mn})\e^{\ii m\varphi}.
\label{eq:posicion}
\end{equation}
Derivando \eqref{eq:general} con respecto a $t$, resulta
\[
\frac{\partial\psi}{\partial t}
=\sum_{m,n}J_{m}\left(\frac{\alpha_{mn}\rho}{a}\right)\e^{\ii m\varphi}
\ii\omega_{mn}\left[A_{mn}\e^{\ii\omega_{mn}t}-B_{mn}\e^{-\ii\omega_{mn}t}\right].
\]
Sustituyendo $t=0$ y usando $\partial\psi/\partial t=\eta$, se obtiene
\begin{equation}
\eta(\rho,\varphi)
=\sum_{m,n}J_{m}\left(\frac{\alpha_{mn}\rho}{a}\right)
\ii\omega_{mn}(A_{mn}-B_{mn})\e^{\ii m\varphi}.
\label{eq:velocidad}
\end{equation}

Para determinar $A_{mn}+B_{mn}$, multiplicamos \eqref{eq:posicion} por $\e^{-\ii q\varphi}$ e integramos en $\varphi$. Usando
\[
\frac{1}{2\pi}\int_0^{2\pi}\e^{\ii(m-q)\varphi}\dd\varphi=\delta_{mq},
\]
queda
\[
\sum_{n=1}^{\infty}J_{m}\left(\frac{\alpha_{mn}\rho}{a}\right)(A_{mn}+B_{mn})
=\frac{1}{2\pi}\int_0^{2\pi}\psi_0(\rho,\varphi)\e^{-\ii m\varphi}\dd\varphi.
\]
Luego multiplicamos por $\rho J_{m}(\alpha_{mn}\rho/a)$ e integramos en $\rho$. Usando la ortogonalidad
\begin{equation}
\int_0^a \rho J_{m}\left(\frac{\alpha_{mn}\rho}{a}\right)
J_{m}\left(\frac{\alpha_{mn'}\rho}{a}\right)\dd\rho
=\frac{a^2}{2}\left[J_{m+1}(\alpha_{mn})\right]^2\delta_{nn'},
\label{eq:ortogonalidad}
\end{equation}
se obtiene
\begin{equation}
A_{mn}+B_{mn}
=\frac{2}{a^2[J_{m+1}(\alpha_{mn})]^2}
\frac{1}{2\pi}\int_0^{2\pi}\int_0^a
\psi_0(\rho,\varphi)J_{m}\left(\frac{\alpha_{mn}\rho}{a}\right)\e^{-\ii m\varphi}\rho\dd\rho\dd\varphi.
\label{eq:suma}
\end{equation}

Del mismo modo, para determinar $A_{mn}-B_{mn}$, multiplicamos \eqref{eq:velocidad} por $\e^{-\ii q\varphi}$, integramos en $\varphi$ y luego usamos \eqref{eq:ortogonalidad}. Así,
\begin{equation}
\ii\omega_{mn}(A_{mn}-B_{mn})
=\frac{2}{a^2[J_{m+1}(\alpha_{mn})]^2}
\frac{1}{2\pi}\int_0^{2\pi}\int_0^a
\eta(\rho,\varphi)J_{m}\left(\frac{\alpha_{mn}\rho}{a}\right)\e^{-\ii m\varphi}\rho\dd\rho\dd\varphi.
\label{eq:resta}
\end{equation}
Definimos
\[
P_{mn}:=\frac{2}{a^2[J_{m+1}(\alpha_{mn})]^2}
\frac{1}{2\pi}\int_0^{2\pi}\int_0^a
\psi_0J_{m}\left(\frac{\alpha_{mn}\rho}{a}\right)\e^{-\ii m\varphi}\rho\dd\rho\dd\varphi,
\]
\[
Q_{mn}:=\frac{2}{a^2[J_{m+1}(\alpha_{mn})]^2}
\frac{1}{2\pi}\int_0^{2\pi}\int_0^a
\eta J_{m}\left(\frac{\alpha_{mn}\rho}{a}\right)\e^{-\ii m\varphi}\rho\dd\rho\dd\varphi.
\]
Entonces \eqref{eq:suma} y \eqref{eq:resta} se escriben como
\[
A_{mn}+B_{mn}=P_{mn},
\qquad
\ii\omega_{mn}(A_{mn}-B_{mn})=Q_{mn}.
\]
Sumando y restando estas ecuaciones, se obtiene finalmente
\begin{equation}
A_{mn}=\frac{1}{2}\left(P_{mn}+\frac{Q_{mn}}{\ii\omega_{mn}}\right),
\qquad
B_{mn}=\frac{1}{2}\left(P_{mn}-\frac{Q_{mn}}{\ii\omega_{mn}}\right).
\label{eq:coeficientes}
\end{equation}

Si la membrana parte desde el reposo, entonces $\eta(\rho,\varphi)=0$. Por la definición de $Q_{mn}$, se tiene $Q_{mn}=0$. Sustituyendo esto en \eqref{eq:coeficientes}, queda
\[
A_{mn}=B_{mn}=\frac{P_{mn}}{2}.
\]
Finalmente, sustituyendo $A_{mn}=B_{mn}=P_{mn}/2$ en \eqref{eq:general}, se obtiene
\begin{equation}
\psi(\rho,\varphi,t)
=\sum_{m=-\infty}^{\infty}\sum_{n=1}^{\infty}
P_{mn}J_{m}\left(\frac{\alpha_{mn}\rho}{a}\right)\e^{\ii m\varphi}
\cos(\omega_{mn}t).
\label{eq:reposo}
\end{equation}
En este caso,
\[
A_{mn}=B_{mn}
=\frac{1}{a^2[J_{m+1}(\alpha_{mn})]^2}
\frac{1}{2\pi}\int_0^{2\pi}\int_0^a
\psi_0(\rho,\varphi)J_{m}\left(\frac{\alpha_{mn}\rho}{a}\right)\e^{-\ii m\varphi}\rho\dd\rho\dd\varphi.
\]
De forma equivalente, también se puede usar una base real de soluciones separables:
\[
J_m\left(\frac{\alpha_{mn}\rho}{a}\right)\cos(m\varphi)\cos(\omega_{mn}t),\quad
J_m\left(\frac{\alpha_{mn}\rho}{a}\right)\cos(m\varphi)\sin(\omega_{mn}t),
\]
\[
J_m\left(\frac{\alpha_{mn}\rho}{a}\right)\sin(m\varphi)\cos(\omega_{mn}t),\quad
J_m\left(\frac{\alpha_{mn}\rho}{a}\right)\sin(m\varphi)\sin(\omega_{mn}t),
\]
con $m=0,1,2,\ldots$ y $n=1,2,3,\ldots$.

\section*{Visualización numérica}

El notebook \href{https://github.com/jrosasep/ayudantias/blob/main/fisica-matematica-2/onda_2d_tambor_bessel/onda_disco_funcion_rara.ipynb}{\texttt{onda\_disco\_funcion\_rara.ipynb}} contiene una implementación numérica de la expansión anterior para una condición inicial arbitraria. En particular, allí se proyecta una función inicial $\psi_0(\rho,\varphi)$ sobre los modos normales de la membrana y se construye una animación de la solución truncada $\psi(\rho,\varphi,t)$.

La figura siguiente muestra un fotograma estático de esa visualización. La animación completa se puede ver en \href{https://github.com/jrosasep/ayudantias/blob/main/fisica-matematica-2/onda_2d_tambor_bessel/onda-disco-funcion-rara-ani.gif}{\texttt{onda-disco-funcion-rara-ani.gif}}.

Esta visualización permite ver cómo la condición inicial se descompone en modos de Bessel y cómo evoluciona la membrana cuando se considera velocidad inicial nula.

\end{document}
